% Template for your submission.pdf. Enter your team name below, write your answers after "Answer:", and compile
% this file to create submission.pdf.
\documentclass[addpoints]{exam}

% ---------- Fill in ----------
\newcommand{\teamname}{my-team}  % your team name, as registered
% ---------- End of the part to fill in ----------

\newcommand{\assignmentnumber}{3}

% ---------- Layout of the assignment sheets (from the sheets' macros.tex) ----------
\usepackage{amsmath, amsfonts, amssymb}
\usepackage{graphicx}
\usepackage[super]{nth}
\usepackage{parskip}
\usepackage[hyperfootnotes=false,colorlinks]{hyperref}
\usepackage[utf8]{inputenc}
\usepackage[T1]{fontenc}
\usepackage{lmodern}
\usepackage[protrusion=false,expansion=false]{microtype}
\DisableLigatures{encoding = T1, family = tt*}
\usepackage{enumitem}
\pagestyle{headandfoot}

\newcommand{\lecture}{Deep Learning}
\newcommand{\semester}{WS 2026/2027}
\newcommand{\assignment}[1]{\nth{#1} Assignment}

\newcommand{\formatqwithpoints}{
\qformat{\large \thequestion. [\thepoints] \textbf{\thequestiontitle}\hfill\vrule depth \spacingbelowquestion width 0pt}
\bonusqformat{\large \thequestion. [\thepoints] \textbf{\thequestiontitle}\hfill\vrule depth \spacingbelowquestion width 0pt}
}
\newcommand{\formatqwithoutpoints}{
\qformat{\large \thequestion. \textbf{\thequestiontitle}\hfill\vrule depth \spacingbelowquestion width 0pt}
\bonusqformat{\large \thequestion. \textbf{\thequestiontitle}\hfill\vrule depth \spacingbelowquestion width 0pt}
}
\newcommand{\spacingbelowquestion}{0.85em}
\formatqwithpoints
\bonuspointpoints{bonus point}{bonus points}
\renewcommand{\thepartno}{\arabic{partno}}
\renewcommand\partlabel{\thepartno)}
\bracketedpoints

\setlength\parindent{0em}
\setlength{\parskip}{1.0em}
\renewcommand{\baselinestretch}{1.0}
\renewcommand{\questionshook}{\setlength{\itemsep}{0.7em}
    \setlength{\leftmargin}{0pt}\setlength{\labelwidth}{-\labelsep}}
\renewcommand{\partshook}{\setlength{\itemsep}{0em}
    \setlength{\leftmargin}{1em}
    \setlength{\labelwidth}{-\labelsep}
}
\setlist[itemize]{itemsep=0em,parsep=0em}
\newcommand{\myhrule}{\rule{\textwidth}{.5pt}}

\newcommand{\titledq}[2]{\titledquestion{#1}[#2]\setlength{\parskip}{1.0em}}
\newcommand{\titledqwithparts}[1]{
    \large \formatqwithoutpoints
    \titledquestion{#1}[]
    \formatqwithpoints
    \setlength{\parskip}{1.0em}
    \normalsize
}
\newcommand{\code}[1]{\texttt{#1}}
\newcommand{\fileref}[1]{\texttt{#1}}
\newcommand{\tbold}[1]{\textbf{#1}}
\newcommand{\hint}{\tbold{Hint:}\ }
\newcommand{\answer}{\tbold{Answer:}\ }

% questions and parts with the number, title and points they have on the sheet
% (the exam class keeps counting 1, 2, 3 internally; only the printed numbers are the sheet's)
\newcommand{\sheetnumber}{}
\newcommand{\sheetpartnumber}{}
\renewcommand{\thequestion}{\sheetnumber}
\renewcommand{\thepartno}{\sheetpartnumber}
\newcommand{\sheetq}[3]{\renewcommand{\sheetnumber}{#1}\titledq{#2}{#3}}
\newcommand{\sheetqwithparts}[2]{\renewcommand{\sheetnumber}{#1}\titledqwithparts{#2}}
\newcommand{\sheetpart}[3]{\renewcommand{\sheetpartnumber}{#1}\part[#2] \tbold{#3}}

\firstpageheader{Team: \teamname}{{\textbf{\lecture}}\\ Submission for the \assignment{\assignmentnumber}}{\semester}
\firstpagefooter{}{\thepage}{}
\runningheader{Team: \teamname}{Submission for the \assignment{\assignmentnumber}}{\semester}
\runningfooter{}{\thepage}{}
\headrule
% ---------- End of the layout ----------

\newcommand{\addtemplatenotes}{
\setlength{\parskip}{0.0em}
\small{
\tbold{About this template:} This template contains the questions from the assignment sheet that require a
written answer, with the original numbering and titles. While you may submit a handwritten scan instead, we
strongly encourage using this \LaTeX{} template. \LaTeX{} is widely used for scientific papers and theses,
especially for mathematical notation, so practicing it here is good preparation for your future research and
thesis.
\begin{itemize}
    \item \tbold{Using this template:} Write your answer after \tbold{Answer:} below each question and compile
    \code{submission.tex} to create \code{submission.pdf}.
    \item \tbold{Showing your work:} Always provide your intermediate steps and reasoning, not just the final
    result. For mathematical questions (derivations, proofs or calculations), show each step. For conceptual
    questions, briefly explain the reasoning behind your answer.
    \item \tbold{Team name \& anonymity:} Enter your team name at the top of \code{submission.tex}. Do not add
    any other personal information, such as GitHub usernames, real names, email addresses or matriculation
    numbers.
    \item \tbold{How to hand in:} Commit and push \code{submission.pdf} to your team's repository together with
    \code{submission.tex} and all files that it uses (for example, images and \code{.bib} files).
\end{itemize}
\myhrule%
}%
\setlength{\parskip}{1.0em}%
}

\begin{document}
    \addtemplatenotes

    \begin{questions}

        \sheetqwithparts{2}{Pen and Paper Tasks}

        Perform a forward and backward pass to calculate the gradients for the weights $w_0$, $w_1$, $w_2$, $w_s$
        in the MLP on the sheet.

        \begin{parts}
            \sheetpart{1}{2}{}Perform the backpropagation algorithm for the network on the sheet.

            First perform the forward pass, followed by the backward pass to obtain gradients for the weights
            $w_0, w_1, w_2, w_s$. Please assign equations/values to variables step by step similar to the lecture,
            and reuse them later when possible to make your solution more readable (see slide 23, Lecture 3).

            \answer Your answer.
            % For math over several lines, uncomment and fill in:
            % \begin{align*}
            %     z_0 &= w_0 x \\
            %     h_0 &= g_0(z_0)
            % \end{align*}

            \sheetpart{2}{1}{}Explain what difference the skip connection makes when propagating back the error.

            \answer Your answer.

            \sheetpart{3}{2}{}Calculate the gradients for the given datapoint and the given initial weights
            (calculating the gradients requires to calculate a forward pass first). Also calculate the weights and the
            loss after one gradient descent step.
            \begin{align*}
                (x_1, y_1) = (1, -3)\\
                w_0 = w_1 = w_2 = w_s = 0.5\\
                LearningRate = 1
            \end{align*}

            \hint Gradient descent is the simplest method to train/optimize your neural network. Given some learning
            rate $\alpha$ and the gradient $\frac{\partial L}{\partial w}$ for a weight $w$, we update the weight like
            this: $w = w - \alpha \frac{\partial L}{\partial w}$.

            \answer Your answer.
        \end{parts}

        \sheetqwithparts{4}{Questions on experiments}

        Answer the following questions based on the experiments of question 3 (\fileref{run\_xor\_experiment.py}):

        \begin{parts}
            \sheetpart{1}{\half}{}What do you observe with regards to the loss after 1 step of updating parameters?

            \answer Your answer.

            \sheetpart{2}{\half}{}What do you observe after multiple steps of updating parameters? Does the loss
            always decrease? Explain why it may not.

            \answer Your answer.

            \sheetpart{3}{\half}{}Run the experiment multiple times. Do you always end up with the correct final
            predictions? Explain why or why not?

            \answer Your answer.

            \sheetpart{4}{\half}{}What is the role of the variable \code{lr} in \fileref{run\_xor\_experiment.py}?
            When would you set it to a relatively larger value and when would you set it to a relatively smaller
            value? Explain.

            \answer Your answer.
        \end{parts}

    \end{questions}
\end{document}
